problem:  
Let \((M^n, g)\) be a compact Einstein manifold admitting a real Killing spinor, so that the metric cone  
\[
(C(M), \bar g = dr^2 + r^2 g)
\]
is Ricci-flat and carries a parallel spinor.  Equivalently, \(C(M)\) has special holonomy \(\mathrm{SU}(m), \mathrm{Sp}(m), G_2,\) or \(\mathrm{Spin}(7)\), depending on the dimension and type of the Killing spinor.  

Consider infinitesimal deformations of the **cone metric** of the form  
\[
\bar g_\epsilon = dr^2 + r^2 (g + \epsilon\, h) + O(\epsilon^2),
\]
where \(h\) is a symmetric 2-tensor on \(M\), and assume that \(\bar g_\epsilon\) remains Ricci-flat to first order. Define \(\mathcal{D}_{\text{cone}}\) to be the linearized space of such deformations modulo diffeomorphisms and scaling.  

Let \(\mathcal{D}_{\text{link}}\) be the space of infinitesimal Einstein deformations of \(g\) on \(M\) that preserve the Killing spinor condition, modulo diffeomorphisms and scaling.  

**Question:**  
Is there a natural linear isomorphism
\[
\mathcal{D}_{\text{link}} \;\cong\; \mathcal{D}_{\text{cone}},
\]
obtained by restricting or lifting deformations via the cone construction?  
If not, can one describe precisely the obstruction map measuring the failure of this correspondence—i.e. the analytic condition relating infinitesimal Ricci-flat cone deformations to infinitesimal Killing-spinor-preserving Einstein deformations on the link?  

What is the geometric meaning of elements in \(\mathrm{Ker}(\mathcal{D}_{\text{cone}}\to\mathcal{D}_{\text{link}})\)—do they correspond to asymptotically conical (AC) Ricci-flat deformations that cannot be represented by a global Einstein deformation on the compact link?  

Why is it a "good" problem:  
This problem distills a precise analytical and geometric relation between **deformation theory of special-holonomy Ricci-flat metrics** and **deformation theory of Einstein metrics with Killing spinors**.  It clarifies the link between global cone geometry and intrinsic spinorial structures on the link, focusing on how infinitesimal Ricci-flat deformations interact with spinorial symmetry constraints.  The question integrates spin geometry, Einstein deformation complexes, and elliptic analysis on conical spaces—bridging two major areas of modern differential geometry that are usually studied separately.  

Although Bär’s cone correspondence is a classical static result, the infinitesimal-to-moduli correspondence of the deformation theories remains largely unexplored beyond specific cases (e.g. Sasaki–Einstein or Calabi–Yau).  Establishing or refuting such an isomorphism would profoundly clarify whether the known Killing-spinor geometries (nearly Kähler, nearly parallel \(G_2\), etc.) are infinitesimally rigid or part of nontrivial moduli governed by the cone’s special holonomy.  Its resolution would yield new insights into both Ricci-flat and Einstein deformation theory, providing a unified analytic framework at the intersection of spin geometry, holonomy, and geometric analysis.